\section{A homological device for ascending tori}
\label{sec:homological}

An ascending torus will have type $\Finfty$ for two distinct reasons.
Factoring its defining endomorphism through a finitely presented group
provides finite presentation.  Two finite systems of commuting
splittings then give the higher-dimensional finiteness: one system
annihilates the relevant homology classes by $2$, and the other
annihilates the same classes by $3$.  The factorization also starts
this second argument, by placing the image of the endomorphism in
a finitely generated subgroup.

We first construct the finite systems of idempotents that index the
splittings.  The resulting criterion is Theorem~\ref{thm:torus}.
Its proof must control a finite set of simplex types uniformly over
all cycles; it will impose no bound on chain lengths or coefficients.

\subsection{Two finite systems of idempotents}

For $\ell\in\{2,3\}$, let $D_\ell$ be the unital $\F_\ell$-algebra
of operators generated by $l_0,l_1,r_0,r_1$ on the free vector space
with basis $\{0,1\}^{\N}$.  The operator $l_i$ prefixes $i$, whereas
$r_i$ deletes an initial $i$ and sends every other basis vector to zero.
Thus
\begin{equation}\label{eq:binary-prefix-relations}
 r_i l_j=\delta_{ij}1,
 \qquad l_0r_0+l_1r_1=1.
\end{equation}
For a finite binary word $v$, write $l_v$ for prefixing $v$, write $r_v$
for deleting $v$, and put $E_v=l_vr_v$.

Every product of these operators is zero or a prefix replacement
$l_vr_w$.  Equality of finite linear combinations of prefix replacements
is decidable.  Indeed, choose $m$ at least as large as every domain
prefix length and expand each term by
\[
 l_vr_w=\sum_{|a|=m-|w|}l_{va}r_{wa}.
\]
Collect terms with the same domain prefix and the same output prefix.
For a fixed domain prefix, distinct output prefixes give distinct
outputs on an input whose remaining tail is not eventually periodic.
To see this, equality of two such outputs with different prefixes
would make the shorter prefix a prefix of the longer and would force
the remaining tail to be periodic.  The collected coefficients therefore
all vanish if and only if the operator vanishes.  This is a finite test
over $\F_\ell$.

Put
\[
 A=E_{00},\quad B'=E_{01},\quad C'=E_1,
 \qquad A_0=E_{000},\quad A_1=E_{001}.
\]
For each of the three prefix-free lists
\[
 (00,01,1),\qquad (000,01,1),\qquad (001,01,1),
\]
the operators $l_vr_w$ form matrix units for a copy of
$M_3(\F_\ell)$ in $D_\ell$.  Only the first of these copies has
identity $1$.  Let $\mathcal P_\ell$ be the union of their sets of
idempotents, together with $I=1$.  This is a finite, explicitly
specified set.  In the following statement, orthogonal means
$EP=PE=0$.

\begin{lemma}\label{lem:idempotents}
Let $T$ be an abelian group, and suppose that elements $z_E\in T$,
indexed by $E\in\mathcal P_\ell$, satisfy
\begin{equation}\label{eq:idempotent-additivity}
 z_{E+P}=z_E+z_P
\end{equation}
whenever $E,P,E+P\in\mathcal P_\ell$ and $E,P$ are orthogonal.
Then $\ell z_I=0$.
\end{lemma}

\begin{proof}
Within a rank-three matrix algebra, we will show that the
$\ell$-multiples attached to its diagonal units agree.  The overlapping
corners then identify these multiples for $A,A_0,A_1,B',C'$,
and the split $A=A_0+A_1$ forces their common value to vanish.

First work in any one of the three matrix algebras.  For a column
vector $u$ and a row covector $\lambda$ with $\lambda u=1$, write
$x(u,\lambda)=z_{u\lambda}$.  Fix $v\ne0$ in $\ker\lambda$.
We claim that
\[
 x(u+v,\lambda)-x(u,\lambda)
\]
is independent of the choice of $u$ with $\lambda u=1$.

Suppose first that $u'-u$ and $v$ are linearly independent.  The vectors
$u,u'-u,v$ then form a basis, so there is a covector $\mu$ with
\[
 \mu u=\mu u'=0,\qquad \mu v=1.
\]
The identity
\[
 u\lambda+v\mu=(u+v)\lambda+v(\mu-\lambda)
\]
expresses the same rank-two idempotent in two ways as a sum of
orthogonal rank-one idempotents.  Applying
\eqref{eq:idempotent-additivity} gives
\begin{equation}\label{eq:rank-one-difference}
 x(u+v,\lambda)-x(u,\lambda)
   =x(v,\mu)-x(v,\mu-\lambda).
\end{equation}
The same calculation for $u'$ gives the same difference.  If instead
$u'-u$ belongs to the span of $v$, choose
$w\in\ker\lambda\setminus\langle v\rangle$ and set $u''=u+w$.
Both $u''-u$ and $u''-u'$ are independent of $v$, so comparison through
$u''$ proves the claim.  This argument also applies over $\F_2$.

Apply the claim successively at $u,u+v,\ldots,u+(\ell-1)v$.
The resulting differences are equal, and their sum is zero.  Hence
\[
 \ell x(u+v,\lambda)=\ell x(u,\lambda).
\]
Transposing the argument shows that the $\ell$-multiple is also
unchanged when the covector is varied with the vector fixed and its
value still equal to one.  For two different basis vectors $e_i,e_j$,
take $\lambda=e_i^*+e_j^*$.  Change the covector at $e_i$, change the
vector from $e_i$ to $e_j$, and then change the covector at $e_j$.
It follows that the $\ell$-multiples of all three diagonal units agree.

The overlaps among our three matrix algebras now give
\[
 \ell z_A=\ell z_{B'}=\ell z_{C'}
            =\ell z_{A_0}=\ell z_{A_1}.
\]
Call their common value $q$.  Since $A=A_0+A_1$, relation
\eqref{eq:idempotent-additivity} gives $q=2q$, and therefore $q=0$.
Finally, $I=(A+B')+C'$, with $A+B'$ an idempotent in the first matrix
algebra.  Two applications of the same relation yield
$\ell z_I=0$.
\end{proof}

\subsection{The ascending-torus theorem}

For a group $U$ and an endomorphism $f:U\to U$, write
\[
 T(U,f)=\langle U,t\mid t^{-1}ut=f(u)\ (u\in U)\rangle.
\]
Here the notation includes all relations already holding in $U$.

\Needspace{18\baselineskip}
\begin{theorem}\label{thm:torus}
Let $f:U\to U$ be an injective endomorphism.  Assume the following.
\begin{enumerate}[label=\textup{(\roman*)}]
\item There are a finitely presented group $P$ and homomorphisms
      $a:U\to P$, $b:P\to U$ with $f=ba$.
\item For each $\ell\in\{2,3\}$ there are homomorphisms
      $h_E:U\to U$, indexed by $E\in\mathcal P_\ell$, such that
      $h_I=f$.  Whenever $E,P',E+P'\in\mathcal P_\ell$ and
      $EP'=P'E=0$, the subgroups $h_E(U)$ and $h_{P'}(U)$ commute
      elementwise, and
      \[
       h_{E+P'}(u)=h_E(u)h_{P'}(u)\qquad(u\in U).
      \]
\end{enumerate}
Then $T(U,f)$ is of type $\Finfty$.
\end{theorem}

We first establish finite presentation.  A product-module criterion
then reduces the remaining assertion to a precise homology-vanishing
statement.  Uniformly controlled fillings will supply that vanishing
through the homology sequence of the ascending torus.

\begin{lemma}\label{lem:torus-presentation}
Under assumption \textup{(i)} of Theorem~\ref{thm:torus}, the group
$T(U,f)$ is finitely presented.  This assertion does not require $f$
to be injective.
\end{lemma}

\begin{proof}
Set $g=ab:P\to P$, and let $s$ be the stable letter of $T(P,g)$.
There is a homomorphism
\[
 T(U,f)\longrightarrow T(P,g),\qquad u\longmapsto a(u),\quad
 t\longmapsto s.
\]
An inverse sends $p\in P$ to $t b(p)t^{-1}$ and sends $s$ to $t$.
The defining conjugation relations verify that these are homomorphisms;
their composites are the identity because
$t f(u)t^{-1}=u$ and $s g(p)s^{-1}=p$.
This is the algebraic interleaving of the two direct systems.  No
injectivity of $g$ is needed.  Finally, a finite presentation of $P$,
together with $s$ and its conjugation relations on a finite generating
set of $P$, is a finite presentation of $T(P,g)$.
\end{proof}

\subsection{A product-module criterion}

The product-module test below is a form of the Bieri--Eckmann
finiteness criterion \citep[Proposition~1.2]{BieriEckmann1974}.
For the cellular passage from homological to geometric finiteness,
see also Brown \citep[Section~3]{Brown1975}.  We give the argument for
the now finitely presented group $J=T(U,f)$: vanishing with products
of regular modules will yield finite skeleta of a classifying space.

All group homology below has integral coefficients, and coefficient
modules are right modules.  For a group $J$ and a set $\mathcal I$, the
notation $(\Z J)^{\mathcal I}$ means the direct product, not the
direct sum, with its coordinatewise right action.

\begin{lemma}\label{lem:product-criterion}
Suppose that $J$ is finitely presented and that
\[
 H_n\bigl(J,(\Z J)^{\mathcal I}\bigr)=0
 \qquad(n\ge2)
\]
for every set $\mathcal I$.  Then $J$ is of type $\Finfty$.
\end{lemma}

\begin{proof}
Start with a finite presentation complex and its simply connected
universal cover.  Inductively, suppose we have a connected CW complex
with fundamental group $J$, finite $n$-skeleton, and
$(n-1)$-connected universal cover, where $n\ge2$.  Its universal-cover
$n$-skeleton will be denoted by $Z$.  Attach cells of dimensions at
least $n+1$, without a finiteness requirement, to obtain a contractible
free $J$-complex $\widetilde Z$.  This can be done successively by
killing homotopy groups.  Its augmented cellular chain complex is a
free resolution over $R=\Z J$, and its modules $C_i$ are finitely
generated for $i\le n$.

Let $\mathcal I$ be the set of all cellular $n$-cycles in $Z$, and put
$M=R^{\mathcal I}$.  For every finite-rank free left $R$-module $C$,
the coordinate map gives an isomorphism
\begin{equation}\label{eq:finite-rank-product-tensor}
 M\otimes_R C\ \longrightarrow\ C^{\mathcal I},
 \qquad (m_i)_{i\in\mathcal I}\otimes c
             \longmapsto (m_i c)_{i\in\mathcal I}.
\end{equation}
Using this in degrees $n$ and $n-1$, the family whose coordinate
indexed by a cycle $z$ is $z$ itself defines a cycle in
$M\otimes_R C_n$.  By hypothesis, it is the boundary of some
$y\in M\otimes_R C_{n+1}$.

Choose an $R$-basis for $C_{n+1}$.  Although that basis can be infinite,
the tensor $y$ uses only finitely many of its elements,
say $c_1,\ldots,c_s$.  Taking coordinates of its boundary equation
shows that every $n$-cycle is an $R$-linear combination of
$\partial c_1,\ldots,\partial c_s$.  Thus the $R$-module of
$n$-cycles in $Z$ is finitely generated.  This argument does not assume
that $C_{n+1}$ has finite rank.

Since $Z$ is $(n-1)$-connected, the natural Hurewicz isomorphism
\citep[Theorem~4.37]{Hatcher2002} identifies
$\pi_n(Z)$ with $H_n(Z;\Z)$.  There are no $(n+1)$-cells in $Z$,
so the latter group is precisely its group of cellular $n$-cycles.
Choose spheres representing the finite list of module generators and
attach their $J$-translates.  Downstairs, this attaches finitely many
$(n+1)$-cells; upstairs, it produces an $n$-connected universal cover.
Induction gives a complex with finitely many cells in each dimension
whose universal cover is weakly contractible.  That cover is a CW
complex, so it is contractible by Whitehead's theorem
\citep[Theorem~4.5]{Hatcher2002}.  The resulting complex is the required
$K(J,1)$.
\end{proof}

For $J=T(U,f)$, finite presentation is already available.  It therefore
remains to prove
\[
 H_n\bigl(J,(\Z J)^{\mathcal I}\bigr)=0
 \qquad(n\ge2)
\]
for every set $\mathcal I$.  A tensor chain with these coefficients
may have arbitrarily large supports in its separate coordinates.
The next argument retains the finite set of simplex types shared by
all coordinates; this is the control needed to fill such chains.

\subsection{Uniformly bounded simplex types}

Let $\mathcal E(U)$ be the simplicial set whose $q$-simplices are
tuples $(u_0,\ldots,u_q)$ of elements of $U$.  Faces delete entries
and degeneracies repeat entries.  Its normalized integral chain
complex is denoted by $C_*(U)$.  Thus tuples with consecutive equal
entries are zero, and the boundary is the alternating sum of faces.
The augmentation sends each vertex to $1$.  Simultaneous left
translation gives a free $U$-action. Its augmented chain complex is
contractible as a complex of abelian groups: prepending the identity
gives a contraction.
Consequently its augmented chain complex is a free resolution of the
trivial left $\Z U$-module.

For a finite symmetric set $A\subseteq U$ containing $1$, let
$\mathcal E_A(U)$ consist of the tuples satisfying
\[
 u_i^{-1}u_j\in A\qquad\text{for all }i,j,
\]
and write $C_*^A(U)$ for its normalized chains.  This is a simplicial
subset.  We call $A$ a \emph{bound}.  A chain operator is controlled
on a given input bound if all its output chains lie in $C_*^{A'}(U)$
for some finite output bound $A'$ in each of the degrees under
consideration.  There is no restriction on the number of terms in an
individual output chain.

For each fixed degree $q$, $C_q^A(U)$ has only finitely many basis
elements up to left translation: translate the first vertex to $1$,
after which every other vertex belongs to $A$.  We refer to these
orbits as simplex types.  Every homomorphism of groups sends a finite
bound to a finite bound.  Unions of finitely many bounds may always
be replaced by one finite symmetric bound containing $1$.

Coning at $1$ fills every reduced cycle in the full complex $C_*(U)$.
It does not, however, preserve a common finite bound: on a simplex
with vertex $u$, it introduces the difference $u$ itself.  This becomes
an obstruction when the same bound must work for all translates and all cycles of
$C_*^A(U)$.  We instead construct fillings after an iterate of $f$,
using bounds independent of the individual cycle.

The first step is a controlled additivity calculation in integral
chain complexes.  Its partial homotopy need not be $U$-equivariant.
Later we assemble these ordinary fillings into fillings with product
coefficients.  We will not tensor the chosen homotopy over $\Z U$.

\begin{lemma}\label{lem:controlled-additivity}
Fix $n\ge1$ and a bound $A$.  Suppose that a power $f^k$ induces a
chain map
$F:C_*^A(U)\to C_*(U)$ admitting a controlled chain homotopy to the
constant map at $1$ through degrees less than $n$.  Let
$h_1,h_2:U\to U$ have elementwise commuting images, and put
$h(u)=h_1(u)h_2(u)$.  There is a finite bound $B$ such that every
$n$-cycle $z\in C_n^A(U)$ satisfies
\[
 [(h f^k)_*z]=[(h_1f^k)_*z]+[(h_2f^k)_*z]
       \quad\text{in }H_n(C_*^B(U)).
\]
The same output bound may be chosen for any fixed finite collection
of such pairs of maps.
\end{lemma}

\begin{proof}
Let $H_q$ be the given homotopy operators for $0\le q<n$, and set
$H_q=0$ in all other nonnegative degrees.  The chain map
\[
 F'=F-\partial H-H\partial
\]
is constant at the identity in degree zero and zero in degrees
$1,\ldots,n-1$.  Since $H_n=0$, it satisfies $F'z=Fz$ for every
$n$-cycle $z$.  Choose a common bound for $F$, $F'$, and the finitely
many operators $H_q$ used here.

The diagonal argument below eliminates mixed product terms using
lower-degree vanishing.  The same mechanism appears in the
diagonal-and-K\"unneth induction for mitotic groups
\citep[Proposition~4.1]{BDH1980}.  Here we need an integral chain
calculation retaining a common finite bound.  We use the classical
Alexander--Whitney and shuffle comparison of product and tensor
chain complexes \citep{EilenbergZilber1953}, recording the formulas
and their support control.
Define the diagonal tensor-chain map by
\begin{equation}\label{eq:aw-diagonal}
 D(u_0,\ldots,u_q)
  =\sum_{i=0}^q
       (u_0,\ldots,u_i)\otimes(u_i,\ldots,u_q).
\end{equation}
The tensor differential has the convention
$\partial(c\otimes d)=\partial c\otimes d+
(-1)^{\deg c}c\otimes\partial d$.
The terms at the two sides of each interior cut cancel, showing that
$D$ is a chain map.  Formula~\eqref{eq:aw-diagonal} descends to
normalized chains.

The shuffle map $S$ sends a tensor
$(u_0,\ldots,u_p)\otimes(v_0,\ldots,v_q)$ to the signed sum of
the chains in $U\times U$ obtained by following all monotone lattice
paths from $(0,0)$ to $(p,q)$ and replacing a lattice point $(i,j)$
by $(u_i,v_j)$.  Its sign is $(-1)^N$, where $N$ counts pairs in which
a vertical step precedes a horizontal step.  Deleting a corner of a
path cancels with the path traversing that corner in the other order;
the remaining faces are exactly the tensor differential.  Hence $S$
is a chain map, also on normalized chains.  Both $D$ and $S$ preserve
the bound in each coordinate.

Let $\Delta$ be the chain map induced by $u\mapsto(u,u)$.
There is a controlled chain homotopy between $\Delta$ and $SD$.
Here is a direct construction.  For a basis simplex $c$ with vertex
set $V$, both maps are supported on the full tuple complex on
$V\times V$.  Construct a homotopy $R$ inductively on degree.  In
degree zero the two maps agree, so take $R_0=0$.  Having defined it
on faces, the chain
\[
 \Delta c-SDc-R\partial c
\]
is a cycle supported on that same tuple complex.  It bounds there by
prepending the vertex $(u_0,u_0)$, where $u_0$ is the first vertex of
$c$.  Choose this cone for $R(c)$ and extend on the normalized basis.
Every coordinate vertex used is a vertex of $c$, so $R$ preserves
the bound in both coordinates.  This proves the needed homotopy
without any assumption about chain lengths.

Write $f^k\times f^k:U\times U\to U\times U$ for the product
homomorphism.  On the input complex,
\[
 \Delta F=(f^k\times f^k)_*\Delta
 \simeq (f^k\times f^k)_*SD=S(F\otimes F)D.
\]
The homotopy is $(f^k\times f^k)_*R$ from the preceding construction;
the last equality is naturality of the shuffle.  Thus it is controlled
on the input bound.  Next, the usual signed tensor operator
\[
 H\otimes F+F'\otimes H
\]
is a homotopy between $F\otimes F$ and $F'\otimes F'$.
In the second summand the value on $c\otimes d$ includes the sign
$(-1)^{\deg c}$.  Indeed its homotopy boundary is
\[
 (F-F')\otimes F+F'\otimes(F-F')
       =F\otimes F-F'\otimes F'.
\]
All chains involved have finite bounds separately in the two
coordinates, because this calculation uses only the already
controlled operators.

For an $n$-cycle $z$, applying $F'\otimes F'$ to $Dz$ leaves only
bidegrees $(0,n)$ and $(n,0)$.  The result is
\[
 (1)\otimes F'z+F'z\otimes(1)
     =(1)\otimes Fz+Fz\otimes(1).
\]
Shuffling a vertex against a chain merely holds that coordinate
fixed.  Finally apply the homomorphism
\[
 U\times U\longrightarrow U,
 \qquad (u,v)\longmapsto h_1(u)h_2(v).
\]
It is a homomorphism because the two images commute elementwise.
For coordinate bounds $B_1,B_2$, a difference in its image belongs
to the finite set $h_1(B_1)h_2(B_2)$.  Thus all preceding homotopies
remain controlled after this map.  Their two surviving terms give
the asserted identity.  Taking a union of output bounds proves the
last assertion.
\end{proof}

\begin{proposition}\label{prop:uniform-fillings}
Under the assumptions of Theorem~\ref{thm:torus}, for every finite
bound $A$ and every $n\ge0$ there are an integer $k\ge0$ and a finite
bound $B$ such that $f^k$ sends every reduced integral $n$-cycle of
$C_*^A(U)$ to a boundary in $C_*^B(U)$.
\end{proposition}

\begin{proof}
The word ``reduced'' matters only in degree zero.  The subgroup
$b(P)$ is finitely generated and contains $f(U)$.  Choose a finite
symmetric generating set for $b(P)$.  Every vertex $f(u)$ can be
joined to $1$ by a finite edge path whose differences belong to that
set.  Consequently every reduced zero-cycle is filled after applying
$f$, with a bound independent of the cycle.  This proves the base
case.  Notice that the paths need not have uniformly bounded lengths.

Suppose that the proposition is known in degrees less than $n$,
where $n\ge1$.  We first obtain the controlled partial homotopy
required by Lemma~\ref{lem:controlled-additivity}.  Begin with the
map induced by $f$ and choose the vertex paths just described for
its degree-zero homotopy to the constant map.  Suppose more generally
that a power of $f$ and its homotopy have been constructed through
degree $q-1$, with $1\le q<n$.  For every normalized $q$-simplex $c$, the
missing boundary equation for its homotopy value has right side
\[
 F(c)-H(\partial c),
\]
where $F$ is the current chain map.  For $q\ge1$ this is a cycle:
apply the already established homotopy equation to $\partial c$.
All these cycles, for all basis simplices $c$, have a common finite
bound.  The inductive assertion in degree $q$ supplies one further
power of $f$ that sends all of them to boundaries at a common bound.
Apply that power also to the previously constructed map and homotopy,
then choose a finite filling for each of these boundary equations.
Extending on the normalized basis defines $H_q$ with the required
control.  For $n=1$ only the initial vertex homotopy is needed.

After finitely many such steps, the map induced by some power $f^k$
has a controlled chain homotopy to the constant map through degrees
less than $n$.
The homotopy is not asserted to be $U$-equivariant.  The same exponent
works on all basis simplices because at each step the induction
assertion applies to every cycle in one bounded complex.

Apply Lemma~\ref{lem:controlled-additivity} to every orthogonal
relation in both finite idempotent systems, enlarging the output
bound once to cover all of them.  For a fixed $n$-cycle $z$, set
\[
 z_E=[(h_Ef^k)_*z]
\]
in the homology of this single bounded output complex.  These classes
satisfy all relations of Lemma~\ref{lem:idempotents}.  The identity
class is the same in both systems, namely
\[
 z_I=[(f^{k+1})_*z].
\]
It is annihilated by both $2$ and $3$, and hence is zero.  The bound
and exponent were chosen before $z$, so the conclusion is uniform
over all cycles of the input bound.  This completes the induction.
\end{proof}

\subsection{Products of regular modules and the HNN sequence}

We now complete the proof of Theorem~\ref{thm:torus}.  Put
$J=T(U,f)$.  Since $f$ is injective, the direct limit of the system
$U\xrightarrow f U\xrightarrow f\cdots$ is an increasing union of
copies of $U$.  Its stage shift is an automorphism, and adjoining
that automorphism gives $J$.  Explicitly, its $i$th copy maps into
$J$ by $u\mapsto t^iut^{-i}$; these maps give the inverse to the
direct-limit description.  In particular $U$ embeds in $J$, and
every element of $J$ has the form $t^p u t^{-q}$ with $p,q\ge0$.

There is a $J$-tree with vertices the left cosets $jU$ and one edge
indexed by $jU$ from $jU$ to $jtU$.  The second endpoint is independent
of the representative because $ut=tf(u)$.  Every vertex has one
outgoing edge, and its height increases by one.  The normal form just
given shows that every forward ray eventually meets the ray from
$U$.  Thus the graph is connected.  An undirected cycle would have a
vertex of least height with two outgoing edges, which is impossible.
It is therefore a tree, with one vertex orbit and one edge orbit,
both having stabilizer $U$.

For any right $\Z J$-module $M$, the tree gives the long exact
homology sequence containing
\begin{equation}\label{eq:ascending-hnn-homology}
 \begin{split}
 H_n(U,M)&\xrightarrow{\,1-F\,}H_n(U,M)
       \longrightarrow H_n(J,M)\\
       &\longrightarrow H_{n-1}(U,M)
          \xrightarrow{\,1-F\,}H_{n-1}(U,M).
 \end{split}
\end{equation}
We specify the chain map and a construction of this sequence.
Induce the homogeneous resolution to obtain
$P_* =\Z J\otimes_{\Z U}C_*(U)$.  On it, the second endpoint map is
\[
 j\otimes c\longmapsto jt\otimes f_*(c).
\]
It is well defined, again because $ut=tf(u)$.  The mapping cone of
the difference of the two endpoint maps resolves $\Z$: first taking
homology of the induced resolutions recovers the augmented cellular
chains of the tree, which are exact.  Tensoring that cone with $M$
gives \eqref{eq:ascending-hnn-homology}.  In the resulting complex
$M\otimes_{\Z U}C_*(U)$, its map $F$ is
\begin{equation}\label{eq:coefficient-shift}
 F(m\otimes c)=mt\otimes f_*(c).
\end{equation}

Take now $M=(\Z J)^{\mathcal I}$ for an arbitrary set $\mathcal I$.
We claim that $F$ is locally nilpotent on $H_i(U,M)$ for every $i>0$.
To interpret chains, let $\mathcal S_q(A)$ be the finite set of
normalized $q$-simplex types with bound $A$.  The freeness of these
types gives
\begin{equation}\label{eq:bounded-product-chains}
 \begin{split}
 M\otimes_{\Z U}C_q^A(U)
 &\cong\bigoplus_{s\in\mathcal S_q(A)} M\\
 &\cong\prod_{\alpha\in\mathcal I}
       \left(\bigoplus_{s\in\mathcal S_q(A)}\Z J\right).
 \end{split}
\end{equation}
The second isomorphism uses finiteness of $\mathcal S_q(A)$.

Every tensor chain involves finitely many simplex types, so lies in
such a bounded complex.  In each coordinate $\alpha$, expand its
$\Z J$ coefficients.  The result is a finite ordinary chain in
tuples contained in left cosets $jU$, with all within-coset differences
in one bound $A$ independent of $\alpha$.  Conversely, any family of
such chains, finite in each coordinate and sharing one bound, defines
a tensor chain by \eqref{eq:bounded-product-chains}.  The number of
terms and the coefficient supports may vary without bound as
$\alpha$ varies.

Let a positive-degree tensor cycle represent a class in $H_i(U,M)$.
In each coordinate its chain splits into finitely many coset chains,
and each is a cycle, because the boundary preserves the coset.
Choose an origin $j$ for each coset.  Formula~\eqref{eq:coefficient-shift}
shows that $F^k$ changes that origin to $jt^k$ and applies $f^k$ to
the relative vertices in $U$.  Proposition~\ref{prop:uniform-fillings}
provides one exponent $k$ and one output bound for all these cycles.
Choose a finite filling in each coordinate and each of its finitely
many cosets.  Their sum remains finite in that coordinate.  Even if
several input cosets have the same output coset, their fillings may
simply be added.  All coordinates still use one common finite set of
output simplex types, so their family is a tensor chain by
\eqref{eq:bounded-product-chains}.  It bounds $F^k$ of the original
cycle.  This proves local nilpotence.

For a locally nilpotent endomorphism, $1-F$ is invertible: its inverse
on any element is the finite geometric sum $1+F+F^2+\cdots$ evaluated
on that element.  The sequence~\eqref{eq:ascending-hnn-homology}
therefore gives
\[
 H_n\bigl(J,(\Z J)^{\mathcal I}\bigr)=0
       \qquad(n\ge2).
\]
No claim about local nilpotence on $H_0$ is needed.  By
Lemma~\ref{lem:torus-presentation}, $J$ is finitely presented.
Lemma~\ref{lem:product-criterion} now proves that it is of type
$\Finfty$, completing the proof of Theorem~\ref{thm:torus}.

\subsection{Finite extensions}

The final passage from pointwise to setwise stabilizers will use the
following closure property.  Its proof reuses the product-module
argument above.

\begin{lemma}\label{lem:finite-extension}
If \(K\) is a normal subgroup of finite index in \(L\) and \(K\) has
type \(\Finfty\), then \(L\) has type \(\Finfty\).
\end{lemma}

\begin{proof}
First \(L\) is finitely presented. Indeed, choose a finite presentation
of \(K\) and representatives of the finite quotient \(L/K\). Add the
representatives as generators, their conjugation relations on the
generators of \(K\), and their multiplication relations modulo \(K\).
These finitely many relations reduce every word to a word in \(K\)
followed by one representative, and give a presentation of \(L\).

Put \(R=\Z L\) and \(A=\Z K\).
The finite-type classifying space for \(K\) gives a free \(A\)-resolution
of \(\Z\) that is finitely generated in each degree. Therefore, for
every index set \(I\),
\[
 H_j\left(K,\prod_I A\right)=0\qquad(j>0):
\]
tensoring that resolution with \(\prod_I A\) gives the product of its
underlying integral chain complexes. Products of exact sequences of
abelian groups are exact.

We construct a degreewise finitely generated free \(R\)-resolution of
\(\Z\). Start with the finite generating set of \(L\), which makes the
augmentation ideal finitely generated. Suppose a partial finite free
resolution has been constructed through degree \(n\geq1\), exact
below that degree. Its restriction to \(A\) is still finite free,
because \([L:K]\) is finite. Complete it arbitrarily to a free
\(R\)-resolution and restrict to \(A\). Index a product of copies of
\(A\) by all cycles in degree \(n\). As in
Lemma~\ref{lem:product-criterion}, this family is one tensor cycle, and
the preceding vanishing supplies a boundary using only finitely many
basis elements \(c_1,\ldots,c_s\) of the next \(A\)-chain module.
For each cycle \(z\), its coordinate equation is
\(z=\sum_{i=1}^s a_i(z)\partial c_i\), with \(a_i(z)\in A\).
Thus these boundaries generate the cycle module over \(A\), hence
also over \(R\); no Noetherian hypothesis is used. This
extends the partial resolution by a finite free \(R\)-module.
Induction completes the construction.

For any product \(\prod_I R\), tensoring this degreewise finite
resolution commutes with the product, so
\(H_j(L,\prod_I R)=0\) for \(j>0\).
The finite presentation already proved and
Lemma~\ref{lem:product-criterion} now give type \(\Finfty\).
\end{proof}

